\documentclass[11pt]{article}
\usepackage{paperreview}
\title{Preparation Memory and Energy-Ordering Reversals\\
in Exact Unforced Flows}
\author{Luqman Mahmood}
\date{1 October 2026}
\hypersetup{pdftitle={Preparation Memory and Energy-Ordering Reversals in Exact Unforced Flows},pdfauthor={Luqman Mahmood}}
\begin{document}
\maketitle
\draftstatus{SM1, v1.0.0}{Analytic and computer-assisted results; specialist novelty review remains open.}
\begin{abstract}
We study phase-dependent energy ordering in two explicit, globally smooth
three-component solutions of the unforced Navier--Stokes equation on the torus.
The solutions are independent of one spatial coordinate: a decaying transverse
shear transports a vertical velocity field. Their Fourier representation is
an infinite dissipative nearest-neighbour chain. Initially identical
mode-energy summaries do not determine subsequent total energy. An exact
cosine law reduces the phase contrast to a two-column Gram entry. We give
an explicit all-time no-reversal estimate for family (i) throughout a nonzero
coupling interval, prove that normalized late amplitudes are entire in coupling, and
derive local stability and parity rules for the number of ordering reversals.
Two computer-assisted examples exhibit one reversal each, in opposite
directions; their certificates include a uniform near-zero truncation bound
and are kept separate from the analytic results. Both preparation energies
decrease throughout. The flows are unforced; no laboratory validation is
involved.
\end{abstract}

\section{Question, baseline and contribution}
A dissipating system can retain information not contained in its initial
energy. The relevant fluid baseline is non-normal shear dynamics: decaying
eigenmodes alone need not determine transient response \cite{Trefethen}.
Our energy-ordering contrast does not require either preparation's total
energy to grow. Relaxation-ordering phenomena, including the Markovian
Mpemba mechanism, offer another comparison \cite{LuRaz}. Here the question
is narrower: for an exact unforced fluid family, can one determine the entire
sign history of the difference between two preparation energies?
Exact decaying bar states and enhanced dissipation for related torus flows
are established baselines \cite{BeckWayne,WeiZhangZhao}. Those results make
clear why the reduction itself is not the proposed novelty. The object here
is the sign history of a cross-preparation energy contrast, with both
energies decreasing, rather than a new shear solution or damping mechanism.
The classical lineage includes Kelvin's sheared disturbances and Orr's
stability analysis \cite{Kelvin,Orr}. Shear-enhanced scalar mixing and
vortex relaxation are also established \cite{RhinesYoung,BernoffLingevitch}.
The reduction belongs to the standard two-and-a-half-dimensional class
\cite{MajdaBertozzi}. In particular, the result is not Orr transient
amplification: neither preparation's total energy grows; the signed
\emph{difference between} their energies changes sign.

The contribution proposed for this manuscript is the explicit family and
its accounted infinite-dimensional sign analysis. Shear reductions,
passive advection--diffusion, phase gauges, Dyson expansions, analytic
perturbation and implicit-function arguments are standard tools.
Specialist review of closely related transient-energy and observability
results remains outstanding; this preprint does not claim completed
external peer review or established priority.

\section{An exact unforced reduction}
On the $2\pi$-periodic three-torus let spatial brackets denote normalized
average, and set $\nu=1/16$. For integer horizontal vectors $k,q_0$ and a
horizontal vector $a$ satisfying
\[
 a\cdot k=0,\qquad a\cdot q_0=2c,
\]
define
\[
 v_\theta(x,t)=a\e^{-\rho t}\cos(k\cdot x+\theta),\quad
 \rho=\nu|k|^2,\qquad
 u_\theta=v_\theta+w_\theta e_z,
\]
where all fields are independent of $z$ and
\[
 \partial_t w_\theta+v_\theta\cdot\nabla w_\theta=\nu\Delta w_\theta,
 \qquad
 w_\theta(x,0)=\sin(q_0\cdot x)+\cos((q_0+k)\cdot x).
\]
The horizontal self-advection vanishes. All terms involving a $z$ derivative
vanish as well, and the remaining convection has zero divergence.
A spatially constant pressure therefore makes $u_\theta$ an exact
unforced incompressible Navier--Stokes solution. Positive diffusion and the
smooth prescribed shear give global smooth evolution for this subclass.

We use the following two realizations:
\begin{center}
\begin{tabular}{c c c c c}
\toprule
Family & $k$ & $q_0$ & $a$ & $\rho$\\
\midrule
(i)&$(8,-1,0)$&$(0,2,0)$&$(c/8,c,0)$&$65/16$\\
(ii)&$(1,1,0)$&$(0,-2,0)$&$(c,-c,0)$&$1/8$\\
\bottomrule
\end{tabular}
\end{center}
The second is a coordinate realization of the chain used in the
computer-assisted examples of Section~6.
For $q_m=q_0+mk$, the chains $q_m$ and $-q_m$ are disjoint.
Writing the positive-chain coefficient as $b_m$ gives
\[
 b_m'=-\lambda_m b_m-\ii c\e^{-\rho t}
    (\e^{\ii\theta}b_{m-1}+\e^{-\ii\theta}b_{m+1}),\qquad
 b(0)=-\frac{\ii}{2}e_0+\frac12e_1.
\]
Here $\lambda_m=\nu|q_m|^2$:
\[
 \lambda_m^{(i)}=\frac{65m^2-4m+4}{16},\qquad
 \lambda_m^{(ii)}=\frac{(m-1)^2+1}{8}.
\]
The vertical energy is $\frac12\langle w^2\rangle=\sum_m|b_m|^2$;
there is no additional factor two. The phase-independent horizontal energy
is $|a|^2\e^{-2\rho t}/4$.

\section{Phase contrast and monotone individual energies}
On real $\ell^2(\Z)$ put $(Jr)_m=r_{m-1}-r_{m+1}$.
Then $J^*=-J$ and $\norm J\leq2$. Let
\begin{equation}\label{eq:chain}
 r_j'=-\Lambda r_j+c\e^{-\rho t}Jr_j,\qquad r_j(0)=e_j,\quad j=0,1.
\end{equation}
A diagonal phase gauge followed by multiplication by $(-\ii)^m$ converts
the Fourier evolution to this real equation.
\begin{proposition}[Exact cosine law]
With $C_c(t)=\ip{r_0(t)}{r_1(t)}$,
\begin{equation}\label{eq:cos}
 E_\theta(t)=A_c(t)-\frac12 C_c(t)\cos\theta,\quad
 A_c=E_v+\frac14(\norm{r_0}^2+\norm{r_1}^2).
\end{equation}
In particular $E_\pi-E_0=C_c$ and
\[
 C_c'(t)=-2\ip{r_0}{\Lambda r_1},\qquad
 C_c(t)=c(\lambda_0-\lambda_1)t^2+O(t^3).
\]
\end{proposition}
\begin{proof}
The transformed initial vector is
$\frac{\ii}{2}(-e_0+\e^{-\ii\theta}e_1)$. Expanding its evolved squared
norm gives \eqref{eq:cos}. In differentiating the real inner product the
two $J$ terms cancel. At zero, $C=C'=0$ and
$C''=2c(\lambda_0-\lambda_1)$. Finite-support initial vectors justify
the derivatives; a convergent time-Taylor series at zero is not assumed.
\end{proof}
For any real or complex column combination $r$,
$(\norm r^2)'=-2\ip r{\Lambda r}\leq0$ (with the Hermitian inner product
for complex combinations). Adding the decreasing horizontal energy proves
that both preparation energies decrease. Ordering reversal is not energy
growth. The initial per-mode Hermitian covariances are independent of
$\theta$, and the corresponding pure Stokes energies do not separate.

\begin{remark}[Prediction from the initial covariances]\label{rem:prediction}
Fix the family, coupling and prediction time. A deterministic energy
predictor given only the initial per-mode Hermitian covariances must return
the same value $z$ for preparations $\theta=0$ and $\theta=\pi$.
The triangle inequality and \eqref{eq:cos} imply
\[
 \max\{|z-E_0|,|z-E_\pi|\}\geq\frac{|C_c|}{2},\qquad
 \frac{(z-E_0)^2+(z-E_\pi)^2}{2}
   =(z-A_c)^2+\frac{C_c^2}{4}.
\]
The midpoint $z=A_c$ attains both two-preparation optima; the second
expression is the squared risk for equal prior probabilities on these
two preparations. This is an elementary consequence of their identical
observations, not a limitation on predictors supplied with distinguishing
phase-sensitive or later-time data. At a crossing $C_c=0$ this particular
prediction floor vanishes, so a crossing need not be the time of greatest
prediction ambiguity.
\end{remark}

\section{A uniform all-time no-reversal band}
\begin{theorem}\label{thm:band}
For family (i), every $0<c\leq1/3$ and every $t>0$ satisfy
\[
 C_c(t)\leq-c\e^{-t/2}
  \min\left\{\frac{37}{1152}t^2,
       \frac{35422319}{10201651200}\right\}<0.
\]
\end{theorem}
\begin{proof}
Let $\Gamma=2c(1-\e^{-\rho t})/\rho$. Contractivity of the heat
semigroup and the ordered integration simplex bound the $n$th Dyson
term by $\e^{-t/4}\Gamma^n/n!$ for a unit column.
Only odd total orders survive the phase-flip contrast. Its first term is
\[
 C_1=-c\left[\e^{-t/2}F_{63/8}(t)
                 -\e^{-65t/8}F_{1/4}(t)\right],
 \quad F_a(t)=\frac{1-\e^{-at}}a .
\]
Thus
\[
 |C-C_1|\leq \e^{-t/2}\{\sinh(2\Gamma)-2\Gamma\}
       \leq \e^{-t/2}R,\quad
 R=\e^{2\Gamma}(2\Gamma)^3/6 .
\]
Set $\delta=61/16$ and $L_1=-\e^{t/2}C_1$.
An integral representation gives $L_1'>0$ and
$L_1\geq c\delta t^2\e^{-(63/8)t}$.
For $0<t\leq1/4$, $\e^2<8$, $2\Gamma\leq4ct$, and
$\e^x\leq(1-x)^{-1}$ give
\[
 \frac{L_1}{ct^2}\geq\frac{61}{128},\qquad
 \frac{R}{ct^2}\leq
 \frac{(32/3)c^2t}{1-4ct}\leq\frac49.
\]
The difference is $37/1152$.
For $t\geq1/4$, monotonicity gives $L_1/c\geq61/2048$ while
\[
 R/c\leq\frac{(64/65)^3}{54(1-64/195)} .
\]
Subtracting yields the stated positive rational constant.
\end{proof}
At $c=0$ the contrast vanishes identically; negative $c$ reverses its sign.
The failure of this sufficient estimate at a larger coupling is not itself
a reversal or a physical phase transition.

\section{Late amplitudes and reversal-count structure}
Let $\lambda_*=\min_m\lambda_m$ and $s$ be its unique minimizing index.
These are $(1/4,0)$ in (i) and $(1/8,1)$ in (ii). Set
$z_j=\e^{\lambda_*t}r_j$ and $H=\Lambda-\lambda_*$.

\begin{theorem}[Analytic late contrast]\label{thm:late}
There are entire functions $\alpha_j(c)$ such that
$z_j(t,c)\to\alpha_j(c)e_s$ locally uniformly in complex $c$. For real $c$,
\[
 \e^{2\lambda_*t}C_c(t)\longrightarrow
 L(c)=\alpha_0(c)\alpha_1(c),\quad
 |\alpha_j^{(n)}(c)|\leq(2/\rho)^n .
\]
The entire function $L$ is odd and nonzero, with
$L'(0)=-8/63$ in (i) and $L'(0)=4$ in (ii).
Its zeros are isolated and finite on each compact real interval.
\end{theorem}
\begin{proof}
On $|c|\leq R$ the bounded-perturbation series is uniformly dominated by
$\sum_n(2R/\rho)^n/n!$. Hence $z_j$ is entire and bounded by
$\e^{2R/\rho}$ uniformly in time.
Let $P$ project onto $e_s$ and let $Q=I-P$. The positive gap $\delta$
of $H$ on $Q$ gives
\[
 \norm{Qz_j(t)}
 \leq \e^{-\delta t}
 +2R\e^{2R/\rho}\int_0^t\e^{-\delta(t-u)}\e^{-\rho u}\dd u\to0
\]
uniformly on the disk. Also
$\norm{(Pz_j)'}\leq2R\e^{2R/\rho}\e^{-\rho t}$ is integrable.
The resulting uniform holomorphic limit is entire. At real $c$, the
normalized propagator is contractive. Differentiating the equation and
integrating its forcing inductively gives
$\norm{\partial_c^n z_j}\leq[2(1-\e^{-\rho t})/\rho]^n$.
Pass to the limit.

For $De_m=(-1)^me_m$, $DHD=H$ and $DJD=-J$, so
$\alpha_j(-c)=(-1)^{s+j}\alpha_j(c)$.
The first Dyson term gives the two stated derivatives of $L$.
The identity theorem now excludes accumulating zeros.
For complex $c$, $L$ is the bilinear product of amplitudes, not an
antiholomorphic squared norm.
\end{proof}

\begin{proposition}[Local stability and parity]\label{prop:count}
Fix $c_0>0$ with $L(c_0)\ne0$, and assume every positive-time zero
of $C_{c_0}$ is simple. Their number is finite and locally constant in $c$.
If $\sigma_0$ and $\sigma_\infty$ are the early and late signs, the number
$n$ of simple crossings satisfies $(-1)^n=\sigma_0\sigma_\infty$.
\end{proposition}
\begin{proof}
The nonzero quadratic early coefficient excludes a common interval next
to zero. Uniform convergence to nonzero $L$ excludes a common late interval.
Simple zeros on the remaining compact time interval cannot accumulate:
an accumulation would have both $C=0$ and $C_t=0$.
The implicit function theorem continues each root, and a positive
off-root compact minimum excludes new ones nearby. Each crossing flips sign.
\end{proof}
Count changes therefore require a finite-time multiple zero or a late-amplitude
cancellation, under the stated hypotheses. This is a necessary event list,
not a global count classification or a claim that every event changes count.

\section{Two certificate-supported reversal examples}
The computations and their analytic supplement \cite{Internal} support:
\begin{center}
\begin{tabular}{c c c c}
\toprule
Family, coupling & Early sign of $C$ & Unique zero bracket & Late sign\\
\midrule
(ii), $c=1/4$ & $+$ & $[2199/128,10997/640]$ & $-$\\
(i), $c=8$ & $-$ & $[1307/2400,109/200]$ & $+$\\
\bottomrule
\end{tabular}
\end{center}
These brackets are nondimensional.

\subsection{Finite-chain error charged to the continuum}
For a chain truncated on $[-M,M]$, $M=24$, let $v_j$ be extended by zero.
With $a_t=\e^{-\lambda_*t}$ and the above $\Gamma$, set
\[
 \beta=a_t\Gamma^M\e^\Gamma/M!,\qquad \eta=2\beta.
\]
Dyson support implies agreement before order $M$, so
$\norm{r_j-v_j}\leq\eta$ and $|C-C_M|\leq2a_t\eta$.
If $p(k)=\max_{|m|\leq k}\lambda_m$, then
\[
 p_i(k)=65k^2/16+k/4+1/4,\quad
 p_{ii}(k)=k^2/8+k/4+1/4 .
\]
Expanding these polynomials in falling factorials bounds the weighted
tail by $\norm{\Lambda(r_j-v_j)}\leq p(M+1)\eta$. Consequently
\[
 |\ip{r_0}{\Lambda r_1}-\ip{v_0}{\Lambda v_1}|
 \leq2p(M)a_t\eta+p(M+1)\eta^2 .
\]
Near zero, a uniform charge is necessary:
\[
 \frac{|C-C_M|}{t^2}
 \leq \frac{4(2c)^{24}\tau^{22}}{24!(1-2c\tau)}
 \quad(0<t\leq\tau<1/(2c)).
\]
The finite-chain certificates control $C_M$; this bound charges the
whole-chain difference uniformly down to $t=0$.
Safe bounds are $C\geq(3/100)t^2$ up to $3/40$ in (ii) and
$-C\geq(299/10)t^2$ up to $1/300$ in (i).

\subsection{Evidence boundary and reproducibility}
The accompanying supplement provides interval sources, comparison inputs,
a replay driver, and the analytic continuum-error and late-time bounds.
The driver starts from unit initial columns and recomputes both finite-chain
trajectories, the zero-window derivative certificates and all 4,224 early
subcells. It checks the sign runs, straddle cells and adjoining interval
endpoints. The zero-window calculation uses the trajectory computed during
the same invocation. These are same-code reproductions, not independent
ODE implementations.

The supplement distinguishes interval predicates from rounded display
endpoints and requires the corrected whole-chain coefficient. Review of
the Cauchy and Dyson bounds is part of checking the continuum result;
matching a sampled trajectory alone is insufficient. Analytic Theorems
\ref{thm:band}--\ref{thm:late} do not depend on the finite calculations.

\section{Prospective use and limitations}
A test of the prediction must fix its preparation, coupling, observation
times, normalization and instrument error before measurement. It must
compare against the same ordinary unforced equations and retain numerical
and measurement errors separately. The cosine law predicts a relational
quantity; a single declining energy record does not reveal the phase.

A finite error budget can hide the contrast early and late while exposing
it at an intermediate time, although the underlying coupling is continuous.
The equality $E_\theta=E_{-\theta}$ has a rigid-motion explanation in the
specified preparation.
Open questions include the global reversal-count diagram, optimal
no-reversal thresholds, and robustness to physical preparation and readout
errors.

\aistatement
\begin{thebibliography}{99}
\bibitem{Kelvin} W. Thomson (Lord Kelvin),
\emph{Stability of fluid motion---Rectilineal motion of viscous fluid between
two parallel planes}, Philos. Mag. \textbf{24}, 188--196 (1887).
\url{https://doi.org/10.1080/14786448708628078}.
\bibitem{Orr} W. M'F. Orr,
\emph{The stability or instability of the steady motions of a perfect liquid
and of a viscous liquid}, Parts I and II, Proc. Roy. Irish Acad. A
\textbf{27}, 9--68 and 69--138 (1907).
\url{https://www.jstor.org/stable/20490590};
\url{https://www.jstor.org/stable/20490591}.
\bibitem{RhinesYoung} P. B. Rhines and W. R. Young,
\emph{How rapidly is a passive scalar mixed within closed streamlines?},
J. Fluid Mech. \textbf{133}, 133--145 (1983).
\url{https://doi.org/10.1017/S0022112083001822}.
\bibitem{BernoffLingevitch} A. J. Bernoff and J. F. Lingevitch,
\emph{Rapid relaxation of an axisymmetric vortex},
Phys. Fluids \textbf{6}, 3717--3723 (1994).
\url{https://doi.org/10.1063/1.868362}.
\bibitem{MajdaBertozzi} A. J. Majda and A. L. Bertozzi,
\emph{Vorticity and Incompressible Flow}, Cambridge University Press (2002).
\url{https://doi.org/10.1017/CBO9780511613203}.
\bibitem{Trefethen} L. N. Trefethen, A. E. Trefethen, S. C. Reddy
and T. A. Driscoll, \emph{Hydrodynamic stability without eigenvalues},
Science \textbf{261}, 578--584 (1993).
\url{https://doi.org/10.1126/science.261.5121.578}.
\bibitem{LuRaz} Z. Lu and O. Raz,
\emph{Anomalous cooling and heating---the Mpemba effect and its inverse},
arXiv:1609.05271 (2016); PNAS \textbf{114}, 5083--5088 (2017).
\url{https://arxiv.org/abs/1609.05271}.
\bibitem{BeckWayne} M. Beck and C. E. Wayne,
\emph{Metastability and rapid convergence to quasi-stationary bar states
for the two-dimensional Navier--Stokes equations},
Proc. Roy. Soc. Edinburgh A \textbf{143}, 905--927 (2013).
\url{https://arxiv.org/abs/1108.3416}.
\bibitem{WeiZhangZhao} D. Wei, Z. Zhang and W. Zhao,
\emph{Linear inviscid damping and enhanced dissipation for the Kolmogorov flow},
Adv. Math. \textbf{362}, 106963 (2020).
\url{https://doi.org/10.1016/j.aim.2019.106963}.
\bibitem{Internal} L. Mahmood, Reproducibility supplement to this paper,
\texttt{supplement/README.md} and \texttt{PROOF.md} (2026).
Distributed with this paper. DOI:
\url{https://doi.org/10.5281/zenodo.23049605}.
\end{thebibliography}
\end{document}
